\documentclass[11pt]{article}

\usepackage[utf8]{inputenc}
\usepackage[T1]{fontenc}
\usepackage{lmodern}
\usepackage[margin=1in]{geometry}
\usepackage{amsmath,amssymb,amsthm}
\usepackage{booktabs}
\usepackage{microtype}
\usepackage[hidelinks]{hyperref}
\usepackage[capitalise]{cleveref}

% Run-in headings italic, never bold.
\makeatletter
\renewcommand\paragraph{\@startsection{paragraph}{4}{\z@}%
  {3.25ex \@plus1ex \@minus.2ex}{-1em}%
  {\normalfont\normalsize\itshape}}
\makeatother

\title{Hierarchical Equal Risk Contribution Is a Corner of a Schur Square,\\
and Two Dials Lead to Minimum Variance}
\author{Peter Cotton\thanks{Microprediction.
Email: \texttt{peter.cotton@microprediction.com}. A companion
verification script is \texttt{paper/verify\_herc\_square.py}
at \texttt{github.com/microprediction/schur}. Every numerical claim in
this paper is checked there in exact rational arithmetic.}}
\date{September 18, 2026}

\newcommand{\R}{\mathbb{R}}
\newcommand{\ones}{\mathbf{1}}
\DeclareMathOperator{\diag}{diag}
\DeclareMathOperator{\Var}{Var}

\newtheorem{proposition}{Proposition}
\newtheorem{lemma}{Lemma}
\newtheorem{corollary}{Corollary}
\newtheorem{remark}{Remark}
\newtheorem{definition}{Definition}

\begin{document}
\maketitle

\begin{abstract}
Hierarchical equal risk contribution cuts a dendrogram into clusters,
holds the inverse-variance portfolio inside each cluster, and budgets
across clusters by inverse cluster variance. We show that its budget rule
is already the one the global minimum-variance portfolio uses, and that
what it lacks is conditioning in two places: each cluster on the other
clusters, and each asset on its cluster mates. Damping the two Schur
complements by $\gamma$ and $\eta$ gives a square with the method at one
corner and the global optimum at the opposite corner. The inner dial has a
closed form that interpolates the naive inverse-variance quantities and
Stevens' regression quantities, keeps every weight positive up to an
explicit frontier, and on an equicorrelated cluster traces the straight
segment from inverse variance to the cluster's minimum-variance portfolio.
Under estimation error the best point can be strictly inside the square,
and we give the two-dial local theorem at the minimum-variance corner, a
sign criterion for the inner dial with both signs realized, and exact
examples, including one whose optimum is $\eta=3/5$ with gaps $1/45$ and
$1/20$ to the two ends.
\end{abstract}

\section{Introduction}

Hierarchical equal risk contribution, HERC \cite{raffinot2018}, is a
two-level rule. A dendrogram is cut into $k$ clusters. Inside each cluster
the weights are inverse variance. Across clusters the budgets follow the
inverse of each cluster's risk. The method is widely implemented
\cite{riskfolio,skfolio,mlfinlab} and widely tested
\cite{trucios2026,huang2022,sjostrand2020}, with results that are mixed and
that no theory has explained.

Block inversion gives the global minimum-variance portfolio the same
two-level shape. Each cluster's block of $\Sigma^{-1}\ones$ is a
minimum-variance direction on the cluster's covariance conditional on
everything else, and the cluster totals are whatever those directions sum
to \cite{cotton2024schur}. Written as a budget rule, the total is the
inverse of the cluster's conditional minimum variance. That is HERC's rule.

So HERC's across-cluster step is not the approximation. The approximations
are that its inside step reads a diagonal, and that neither step reads
anything outside the cluster. Both are the omission of a Schur complement.
Damping the two complements by separate parameters gives a square. HERC
sits at $(\gamma,\eta)=(0,0)$ and the global minimum-variance portfolio at
$(1,1)$.

The inner dial is new and simple. A single asset conditioned on its cluster
mates is Stevens' regression identity \cite{stevens1998}, and damping it
gives a weight whose numerator and denominator are each linear in $\eta$
between the naive quantity and the regression quantity. The whole inner
dial costs one inverse per cluster. The outer dial is the flat bridge of
\cite{cotton2026nco}, with the outer optimizer replaced by HERC's inverse
fitness.

\Cref{sec:where} asks where to sit when the covariance is estimated. The
inner dial on an equicorrelated cluster reduces to one scalar, which gives
exact interior optima and a sign criterion at the full-coupling end. At the
minimum-variance corner the two-dial local theorem says which dials move
inside and by how much, and exact examples show both dials moving, and one
moving while the other stays.

Everything exact here concerns the unconstrained, fully invested,
long-short problem, as in \cite{cotton2024schur,cotton2026bridge}. HERC is
long-only by construction, and \cref{prop:frontier} says exactly how far
the inner dial can be turned before that is lost.

\section{The method}

Let $\Sigma\succ0$ be the covariance of $n$ assets and let $C_1,\dots,C_k$
be the clusters at which HERC cuts its dendrogram. For an index set $C$
write $\Sigma_{CC}$ for its block, $-C$ for its complement, and $\sigma_i^2$
for $\Sigma_{ii}$.

HERC holds inside cluster $C$ the inverse-variance portfolio
$v_C\propto\diag(\Sigma_{CC})^{-1}\ones$, normalized to $\ones^\top v_C=1$,
and measures the cluster's risk by that portfolio's variance,
$\nu_C=v_C^\top\Sigma_{CC}v_C$. It then budgets across clusters top down
along the dendrogram.

\paragraph{The split rule.} Raffinot writes the share of the first child at
a split as $\mathrm{RC}_1/(\mathrm{RC}_1+\mathrm{RC}_2)$, ``where
$\mathrm{RC}_1$ is the risk contribution of the first cluster''
\cite{raffinot2018}, and does not define $\mathrm{RC}$. Read as a risk, the
formula gives the riskier child more. Read as an inverse risk, a precision,
it gives the riskier child less, which is the intent, and Palomar's
textbook makes the same reading in recovering the inverse-variance split
\cite{palomar2025}. We take $\mathrm{RC}$ of a child to be the sum of the
precisions $1/\nu_C$ of the final clusters it contains.

\begin{proposition}[The dendrogram telescopes]
\label{prop:telescope}
With additive precisions, top-down division along any dendrogram gives
final cluster $C$ the budget $a_C\propto1/\nu_C$. The tree affects the
allocation only through the choice of clusters.
\end{proposition}

\begin{proof}
At a node with children $L$ and $R$ and precisions $P_L=\sum_{C\subseteq L}1/\nu_C$
and $P_R$, the child $L$ receives the fraction $P_L/(P_L+P_R)$ of the node's
budget. Along the path from the root to $C$ the fractions telescope to
$(1/\nu_C)/P_{\mathrm{root}}$.
\end{proof}

So HERC with variance is the flat two-level rule
\begin{equation}
  w \;\propto\; \sum_{C}\frac{v_C}{\nu_C},
  \qquad v_C\propto\diag(\Sigma_{CC})^{-1}\ones,\quad
  \nu_C=v_C^\top\Sigma_{CC}v_C .
  \label{eq:herc}
\end{equation}
In words: hold inverse variance inside each cluster, and give each cluster
capital in proportion to the inverse variance of what it holds.

\begin{remark}[The reference implementations]
The libraries \cite{riskfolio,skfolio,mlfinlab} implement the split as one
minus the ratio of the children's risks, with the risk of a child taken as
the sum of the variances $\nu_C$ of its final clusters. For two clusters
this agrees with \eqref{eq:herc}. For three or more it does not, because
sums of variances do not telescope, and the allocation then depends on the
shape of the dendrogram above the cut. The verification script exhibits
both facts. We analyze \eqref{eq:herc}.
\end{remark}

\section{Block inversion}

Fix a vector $u$ and a cluster $C$. Block inversion of $\Sigma$ on the
partition $(C,-C)$ gives \cite[App.~A]{cotton2024schur}
\begin{equation}
  \big(\Sigma^{-1}u\big)_{C}
  \;=\; \big(\Sigma^{c}_{C}\big)^{-1} b_C,
  \qquad
  \Sigma^{c}_{C} = \Sigma_{CC}-\Sigma_{C,-C}\Sigma_{-C,-C}^{-1}\Sigma_{-C,C},
  \quad
  b_C = u_{C}-\Sigma_{C,-C}\Sigma_{-C,-C}^{-1}u_{-C} .
  \label{eq:blockinv}
\end{equation}
The complement is the covariance of $C$ conditional on the other assets and
$b_C$ is the companion vector. For positive definite $Q$ and a vector $b$
the constrained problem of minimizing $x^\top Qx$ subject to $b^\top x=1$
has solution and attained variance
\begin{equation}
  v(Q,b)=\frac{Q^{-1}b}{b^\top Q^{-1}b},
  \qquad
  \nu(Q,b)=\frac{1}{b^\top Q^{-1}b},
  \qquad\text{so}\qquad
  Q^{-1}b=\frac{v(Q,b)}{\nu(Q,b)} .
  \label{eq:fitness}
\end{equation}
The unnormalized direction is the normalized portfolio divided by its
variance.

Apply this to \eqref{eq:blockinv} with $u=\ones$. The global
minimum-variance portfolio is $\sum_C v_C/\nu_C$ with
$v_C=v(\Sigma^c_C,b_C)$ and $\nu_C=\nu(\Sigma^c_C,b_C)$, which is
\eqref{eq:herc} with two substitutions undone. HERC replaces the
conditional pair $(\Sigma^c_C,b_C)$ by the raw pair $(\Sigma_{CC},\ones)$,
and it replaces the minimum-variance portfolio of that pair by the
inverse-variance portfolio. The budget rule is untouched.

Noguer i Alonso lists three substitutions that turn the exact bisection
identity into hierarchical risk parity \cite{lopezdeprado2016,noguer2026}. HERC makes the
first two, the raw block for the complement and $\ones$ for $b$, together
with the diagonal for the block inside each cluster. It does not make the
third, because its inter-cluster budget is the exact one.

\paragraph{One asset at a time.} Take $C=\{i\}$ in \eqref{eq:blockinv} with
$u=\ones$. Then $\Sigma^c_{\{i\}}$ is the residual variance of asset $i$
regressed on all the others and $b_{\{i\}}$ is one minus the sum of the
regression coefficients. This is Stevens' identity \cite{stevens1998},
\begin{equation}
  \big(\Sigma^{-1}\ones\big)_i
  \;=\;\frac{1-\sum_{j\ne i}\beta_{ij}}{\sigma_i^2\,(1-R_i^2)} ,
  \label{eq:stevens}
\end{equation}
with $R_i^2$ the multiple correlation and $\beta_{ij}$ the hedge
coefficients. The inverse-variance portfolio is the same expression with
the hedges set to zero. Inside a cluster, HERC is therefore Stevens with
the regression switched off.

\section{The square}
\label{sec:square}

Two Schur complements are missing from \eqref{eq:herc}: that of a cluster
against the other clusters, and that of an asset against its cluster
mates. Damp each.

\begin{definition}[The two dials]
For $\gamma\in[0,1]$ the \emph{cluster pair} of $C$ is
\begin{equation}
  Q_C(\gamma)=\Sigma_{CC}-\gamma\,\Sigma_{C,-C}\Sigma_{-C,-C}^{-1}\Sigma_{-C,C},
  \qquad
  b_C(\gamma)=u_C-\gamma\,\Sigma_{C,-C}\Sigma_{-C,-C}^{-1}u_{-C} .
  \label{eq:outer}
\end{equation}
For $\eta\in[0,1]$ and a pair $(Q,b)$ on $C$, the \emph{leaf direction}
$z(Q,b,\eta)$ has entries
\begin{equation}
  z_i=\frac{b_i-\eta\,Q_{i,-i}Q_{-i,-i}^{-1}b_{-i}}
           {Q_{ii}-\eta\,Q_{i,-i}Q_{-i,-i}^{-1}Q_{-i,i}},
  \qquad i\in C,
  \label{eq:inner}
\end{equation}
where $-i$ now means $C\setminus\{i\}$. The portfolio at $(\gamma,\eta)$ is
\begin{equation}
  w_{\gamma\eta}\;\propto\;\sum_C\frac{v_C}{\nu_C},
  \qquad
  v_C=\frac{z_C}{b_C^\top z_C},\quad
  \nu_C=v_C^\top Q_C v_C,\quad
  z_C=z\big(Q_C(\gamma),b_C(\gamma),\eta\big).
  \label{eq:square}
\end{equation}
\end{definition}

The dial $\gamma$ conditions a cluster on the outside world. The dial
$\eta$ conditions an asset on its cluster mates. The budget rule is HERC's
own, the inverse variance of what the cluster holds, measured on the
cluster's pair. We take $u=\ones$ unless said otherwise.

\begin{proposition}[The four corners]
\label{prop:corners}
\leavevmode
\begin{enumerate}
\item $(0,0)$ is HERC, \eqref{eq:herc}.
\item $(1,1)$ is the global direction $\Sigma^{-1}u$, so with $u=\ones$
  the global minimum-variance portfolio.
\item $(0,1)$ holds the minimum-variance portfolio of each cluster's own
  block and budgets by inverse cluster minimum variance.
\item $(1,0)$ holds the inverse-variance portfolio of each cluster's
  conditional pair, with the companion vector in the numerator, and budgets
  by its inverse variance on the pair.
\end{enumerate}
\end{proposition}

\begin{proof}
At $\eta=0$ the leaf direction is $z_i=b_i/Q_{ii}$, which at $\gamma=0$ is
$1/\sigma_i^2$. At $\eta=1$ the leaf direction is the block inversion
\eqref{eq:blockinv} of the pair $(Q,b)$ on the partition
$(\{i\},C\setminus\{i\})$, so $z=Q^{-1}b$, and by \eqref{eq:fitness}
$v_C/\nu_C=Q^{-1}b$. At $\gamma=1$ the pair is $(\Sigma^c_C,b_C)$ and
$Q^{-1}b=(\Sigma^{-1}u)_C$ by \eqref{eq:blockinv}. The remaining corners
read off the definitions.
\end{proof}

\begin{table}[ht]
\centering
\begin{tabular}{@{}lll@{}}
\toprule
corner & inside a cluster & across clusters \\
\midrule
$(0,0)$ HERC & inverse variance on $\Sigma_{CC}$ & inverse variance of the holding \\
$(0,1)$ & minimum variance on $\Sigma_{CC}$ & inverse minimum variance \\
$(1,0)$ & inverse variance on the conditional pair & inverse variance of the holding \\
$(1,1)$ global & minimum variance on the conditional pair & inverse conditional minimum variance \\
\bottomrule
\end{tabular}
\caption{The four corners with $u=\ones$. The across-cluster rule is the
same at every corner; only what it measures changes.}
\label{tab:corners}
\end{table}

The corner $(0,1)$ is a cluster-then-optimize rule with inverse-variance
budgets, of the kind used in \cite{arevalo2019} with a different inner
objective. The
corner $(1,0)$ has no name. Along the edge $\eta=1$ the square is the flat
bridge of \cite{cotton2026nco} with the outer optimizer replaced by
HERC's inverse fitness; the two agree at $\gamma=1$ and differ inside.

\begin{remark}[Cheaper conditioning sets]
The pair \eqref{eq:outer} conditions on all $n-|C|$ other assets, which is
the inverse a two-level method exists to avoid. Two devices from earlier
work apply unchanged. Conditioning composed top down along the dendrogram,
with each block conditioned on its sibling as in \cite{cotton2024schur},
reaches the same pair at $\gamma=1$ because Schur complements compose, and
the script checks this on an unbalanced tree. Conditioning on one knot per
other cluster, as in \cite{cotton2026nco}, is exact under the gateway model
stated there. The tree device inverts sibling blocks, and the knot device
inverts only cluster blocks and one $k\times k$ matrix.
\end{remark}

\begin{remark}[Other companion vectors]
\label{rem:companion}
\Cref{prop:corners} holds for any $u$. With $u=\sigma$, the vector of
volatilities, the corner $(1,1)$ is the most diversified portfolio
$\Sigma^{-1}\sigma$ \cite{choueifaty2008}, and the corner $(0,0)$ holds the
inverse-volatility portfolio inside each cluster with the cash budget of
cluster $C$ proportional to $\mathrm{DR}_C^2/\bar\sigma_C$, where
$\mathrm{DR}_C$ is the diversification ratio of the holding and
$\bar\sigma_C$ its average volatility. HERC with volatility as the risk
measure, as implemented in \cite{riskfolio}, holds the same inside
portfolio but budgets by $\mathrm{DR}_C/\bar\sigma_C$. It is not a corner
of this square, and the diversification ratio enters squared in the rule
that is. With $u=\mu$ the corner $(1,1)$ is the tangency direction, and
$(0,0)$ holds $\mu_i/\sigma_i^2$ inside each cluster with budgets
proportional to the squared Sharpe ratio of the holding divided by its
mean.
\end{remark}

\section{The inner dial}
\label{sec:inner}

The leaf direction \eqref{eq:inner} looks like $|C|$ solves of size
$|C|-1$. It is one solve.

\begin{proposition}[Closed form]
\label{prop:closed}
Let $q_i=1/(Q^{-1})_{ii}$ and $s_i=(Q^{-1}b)_i/(Q^{-1})_{ii}$. Then
\begin{equation}
  z_i(\eta)=\frac{(1-\eta)\,b_i+\eta\,s_i}{(1-\eta)\,Q_{ii}+\eta\,q_i} .
  \label{eq:closed}
\end{equation}
\end{proposition}

\begin{proof}
Block inversion on $(\{i\},-i)$ gives
$Q_{ii}-Q_{i,-i}Q_{-i,-i}^{-1}Q_{-i,i}=1/(Q^{-1})_{ii}$ and
$b_i-Q_{i,-i}Q_{-i,-i}^{-1}b_{-i}=(Q^{-1}b)_i/(Q^{-1})_{ii}$. Both
expressions in \eqref{eq:inner} are affine in $\eta$ between their values
at $0$ and $1$.
\end{proof}

At $\eta=1$ the weight is $(Q^{-1}b)_i$, and with $b=\ones$ the numbers
$q_i$ and $s_i$ are Stevens' residual variance and one minus the hedge sum
\eqref{eq:stevens}. The dial interpolates both the numerator and the
denominator of Stevens' formula linearly. The denominator is a convex
combination of two positive numbers, so it never vanishes.

\begin{proposition}[Long-only frontier]
\label{prop:frontier}
Let $b>0$ entrywise. The leaf direction $z(\eta)$ is nonnegative in every
entry if and only if $\eta\le\eta_+$, where
\begin{equation}
  \eta_+=\min_{i:\,s_i<0}\frac{b_i}{b_i-s_i},
  \label{eq:frontier}
\end{equation}
and $\eta_+=1$ when no $s_i$ is negative. The first short appears at the
asset whose hedge sum first exceeds its companion entry.
\end{proposition}

\begin{proof}
The denominator in \eqref{eq:closed} is positive. The numerator is affine
in $\eta$, equals $b_i>0$ at $\eta=0$, and is negative at $\eta=1$ exactly
when $s_i<0$, crossing zero at $\eta=b_i/(b_i-s_i)$.
\end{proof}

HERC is long-only by construction and the global minimum-variance portfolio
is not. The frontier says how much of the inside conditioning can be
restored before that changes, cluster by cluster, at the cost of the
inverse already computed.

\paragraph{The equicorrelated cluster.} Let cluster $C$ have $m$ assets
with volatilities $\sigma_i$ and a common correlation $\rho$, and let
$\gamma=0$ so the pair is $(\Sigma_{CC},\ones)$. Write
$\kappa=\rho/(1+(m-2)\rho)$ for the regression coefficient of one
standardized asset on another, and $s_{-i}=\sum_{j\ne i}1/\sigma_j$.

\begin{proposition}[A straight segment]
\label{prop:segment}
On an equicorrelated cluster,
\begin{equation}
  z_i(\eta)\;\propto\;\frac{1-\eta\,\kappa\,\sigma_i s_{-i}}{\sigma_i^2},
  \label{eq:equi}
\end{equation}
with a factor common to all $i$. The normalized holding lies on the
straight segment from the inverse-variance portfolio $v_0$ to the
cluster's minimum-variance portfolio $v_\star$,
\begin{equation}
  v(\eta)=v_\star+\big(1-\theta(\eta\kappa)\big)(v_0-v_\star),
  \qquad
  \theta(t)=\frac{t\,(\alpha-\kappa\beta)}{\kappa\,(\alpha-t\beta)},
  \label{eq:theta}
\end{equation}
where $\alpha=\sum_i\sigma_i^{-2}$ and $\beta=\sum_i s_{-i}/\sigma_i$, so
that $\beta=(\sum_i\sigma_i^{-1})^2-\alpha$. The long-only frontier is
$\eta_+=\min_i 1/(\kappa\sigma_i s_{-i})$.
\end{proposition}

\begin{proof}
In the equicorrelated model every asset has the same multiple correlation
on the rest, so the denominators in \eqref{eq:inner} share one factor, and
the hedge sum of asset $i$ is $\kappa\sigma_i s_{-i}$. So the unnormalized
weight is $a_i-t\,c_i$ with $t=\eta\kappa$, $a_i=\sigma_i^{-2}$ and
$c_i=s_{-i}/\sigma_i$. Two normalized vectors of this form differ by a
multiple of the fixed vector $a\beta-c\alpha$, which gives the segment,
and $t=\kappa$ is the minimum-variance end. The frontier is
\eqref{eq:frontier} with $b=\ones$ and $s_i=1-\kappa\sigma_i s_{-i}$.
\end{proof}

In words: on an equicorrelated cluster, the inner dial is linear shrinkage
between the two familiar portfolios, with the mixing weight
\eqref{eq:theta} a fixed function of the estimated correlation. With equal
volatilities the two ends coincide and the dial does nothing.

\paragraph{The cut.} With one cluster, $\gamma$ is idle and the square is
the inner dial applied to all of $\Sigma$, from inverse variance to global
minimum variance. With $n$ singleton clusters, $\eta$ is idle and the pair
of $\{i\}$ is Stevens' pair, so the outer dial traces the same path. HERC
is the inverse-variance portfolio at both extremes of the cut, and departs
from it only through the budgets \eqref{eq:herc} at intermediate cuts.
The dendrogram chooses where on that path the two dials separate.

\section{Where to sit}
\label{sec:where}

Let $\widehat\Sigma_\tau=\Sigma+\tau E$ be an estimate, where $E$ is a
random symmetric matrix whose law is unchanged by $E\mapsto-E$, and
$\widehat\Sigma_\tau\succ0$ on the range of $\tau$ considered. The clusters
are held fixed. Every quantity in \eqref{eq:square} is computed from
$\widehat\Sigma_\tau$ and normalized by $\ones^\top w=1$. Define
\begin{equation}
  F(\gamma,\eta,\tau)=\mathbb{E}\big[w_{\gamma\eta}(\widehat\Sigma_\tau)^\top
  \,\Sigma\,w_{\gamma\eta}(\widehat\Sigma_\tau)\big],
  \qquad V_0=F(\cdot,\cdot,0),
  \qquad G=\tfrac12\,\partial_\tau^2F(\cdot,\cdot,0),
  \label{eq:F}
\end{equation}
and assume the expansion $F=V_0+\tau^2G+O(\tau^4)$ holds in $C^2$ in the
dials, as in \cite{cotton2026nco}.

\begin{lemma}[The HERC edge is blind]
\label{lem:blind}
For every $\eta$, the portfolio $w_{0\eta}$ depends on
$\widehat\Sigma_\tau$ only through its diagonal blocks
$\widehat\Sigma_{CC}$.
\end{lemma}

\begin{proof}
At $\gamma=0$ the pair of $C$ is $(\widehat\Sigma_{CC},\ones)$, the leaf
direction and the fitness read only the pair, and the budget rule reads
only the fitnesses.
\end{proof}

So under noise supported off the diagonal blocks, $F(0,\eta,\tau)=V_0(0,\eta)$
exactly, and the argument of \cite{cotton2026bridge} transfers: the cost of
turning $\gamma$ away from zero is of order $\gamma\tau^2$, and at any
$\eta$ with $\partial_\gamma V_0(0,\eta)<0$ the point $(0,\eta)$ is
strictly suboptimal for small $\tau$. Under whole-matrix noise the edge is itself
noisy and the comparison is between noise costs, as discussed there.

\begin{proposition}[The minimum-variance corner, two dials]
\label{prop:corner}
Let $J=[\partial_\gamma w,\ \partial_\eta w]$ at $(1,1,0)$ and suppose its
columns are linearly independent. Then $\nabla V_0(1,1)=0$, the Hessian is
$H=2J^\top\Sigma J\succ0$, and for all sufficiently small $\tau$ the map
$F(\cdot,\cdot,\tau)$ has a unique stationary point near $(1,1)$, a strict
local minimizer, at
\begin{equation}
  (\tilde\gamma,\tilde\eta)=(1,1)-H^{-1}\nabla G(1,1)\,\tau^2+O(\tau^4).
  \label{eq:shift}
\end{equation}
If both components of $H^{-1}\nabla G(1,1)$ are positive the local
minimizer is strictly inside the square for small $\tau>0$. If one is
negative the stationary point lies outside the square, the constrained
minimizer lies on one of the two edges through the corner, and its position
along an edge is given by the one-dial formula with the edge's own second
derivatives.
\end{proposition}

\begin{proof}
Since $\ones^\top w=1$, every derivative of $w$ in the dials has zero sum.
By \cref{prop:corners}, $\Sigma w_{11}\propto\ones$, so first derivatives
of $V_0$ vanish and in the second derivatives the terms containing second
derivatives of $w$ vanish, leaving $2J^\top\Sigma J$. The implicit function
theorem applied to $\nabla_{\gamma\eta}F=0$ at $(1,1,0)$ gives
\eqref{eq:shift} and strict convexity on a neighbourhood not depending on
$\tau$. The edge statement is the one-dial case of the same argument.
\end{proof}

The displacement is of second order in the noise and the gain from making
it is of fourth order, as in \cite{cotton2026nco}. The structure of the
square does not fix the signs; both occur below.

\subsection{The inner dial, exactly}

Let $\gamma=0$ and take a single equicorrelated cluster as in
\cref{prop:segment}, with the volatilities known and the correlation
estimated, $\widehat\rho=\rho+\tau Z$ for a symmetric $Z$. Then the recipe
holds $v(\eta)$ at $t=\eta\,\kappa(\widehat\rho)$, while the true portfolio
that minimizes variance is at $t=\kappa(\rho)$. By \eqref{eq:theta} and the
segment,
\begin{equation}
  F(\eta,\tau)=V_\star+\Delta\;\mathbb{E}\Big[\big(1-\theta(\eta\,\kappa(\widehat\rho))\big)^2\Big],
  \qquad
  1-\theta(t)=\frac{(\kappa-t)\,\alpha}{\kappa\,(\alpha-t\beta)},
  \label{eq:Fleaf}
\end{equation}
where $V_\star$ is the cluster's minimum variance and
$\Delta=V(v_0)-V_\star$ is the excess variance of inverse variance over it.
The inner dial's noise problem is one scalar, and nothing else about the
cluster enters.

\paragraph{An exact interior optimum.} Let the estimate be
$\widehat\rho\in\{0,2\rho\}$ with equal probability. On the branch
$\widehat\rho=0$ the dial does nothing. On the branch $\widehat\rho=2\rho$
the choice
\begin{equation}
  \eta_\star=\frac{\kappa(\rho)}{\kappa(2\rho)}
  =\frac{1+2(m-2)\rho}{2\,\big(1+(m-2)\rho\big)}
  \label{eq:etastar}
\end{equation}
puts $t$ at $\kappa(\rho)$ and hence $\theta=1$, the population optimum.
One branch is constant and the other attains its minimum, uniquely because
$\theta$ is strictly monotone in $t$. So $\eta_\star$ is the unique global
optimum, strictly inside, independent of the volatilities. With three
assets of volatilities $(1,2,2)$ and $\rho=\tfrac14$ the optimum is
$\eta_\star=\tfrac35$, and in exact arithmetic
$F(0)-F(\tfrac35)=\tfrac1{45}$ and $F(1)-F(\tfrac35)=\tfrac1{20}$.

\begin{proposition}[Sign at the full-coupling end]
\label{prop:sign}
Let $\widehat\rho=\rho\pm\tau$ with equal probability, $q=(m-2)\rho$, and
$r=\beta/\alpha$. Then $\partial_\eta F(1,\tau)$ has the sign of
\begin{equation}
  \alpha\,(1-q)+\kappa\beta\,(2+q),
  \qquad\text{equivalently}\qquad
  q^2-1<r\,\rho\,(2+q),
  \label{eq:crit}
\end{equation}
for all sufficiently small $\tau$, and the optimum is interior when this
holds and stays at $\eta=1$ when the reverse strict inequality holds. The
interior optimum sits at
\begin{equation}
  \tilde\eta=1-\frac{\alpha(1-q)+\kappa\beta(2+q)}{\rho^2(1+q)^2(\alpha-\kappa\beta)}\,\tau^2+O(\tau^4).
  \label{eq:etatilde}
\end{equation}
These statements are global on $[0,1]$, since $V_0$ has the unique
minimizer $\eta=1$ and $F(\cdot,\tau)\to V_0$ uniformly.
\end{proposition}

\begin{proof}
Differentiate \eqref{eq:Fleaf} at $\eta=1$:
$\partial_\eta F(1,\tau)=2\Delta\alpha^2(\alpha-\kappa\beta)\kappa^{-2}\,
\mathbb{E}[\varphi(\kappa(\widehat\rho))]$ with
$\varphi(x)=x(x-\kappa)/(\alpha-x\beta)^3$. Expand
$\mathbb{E}[\varphi(\kappa(\rho\pm\tau))]$ to second order using
$\varphi(\kappa)=0$, $\varphi'(\kappa)=\kappa/(\alpha-\kappa\beta)^3$,
$\varphi''(\kappa)=2(\alpha+2\kappa\beta)/(\alpha-\kappa\beta)^4$,
$\kappa'=(1+q)^{-2}$ and $\kappa\kappa''=-2q(1+q)^{-4}$. The bracket
collects to $[\alpha(1-q)+\kappa\beta(2+q)]/(1+q)^4$ over
$(\alpha-\kappa\beta)^4$. Dividing by $V_0''(1)=2\Delta\alpha^2/(\alpha-\kappa\beta)^2$
gives \eqref{eq:etatilde}. Multiplying \eqref{eq:crit} by
$(1+q)/\alpha$ gives the second form.
\end{proof}

In words: the dial should stop short of full coupling unless the cluster is
both large and highly correlated, with $(m-2)\rho$ well above one, and its
precision is concentrated in one asset, so that $r$ is small. Equal
volatilities give $r=m-1$, which always satisfies \eqref{eq:crit}, but there
the two ends coincide. The script confirms both signs: volatilities
$(1,2,3,5)$ with $\rho=\tfrac1{10}$ move inside, and volatilities
$(1,20,30,50)$ with $\rho=\tfrac9{10}$ stay at full coupling.

\paragraph{Why the sign can go either way.} The mixing weight
$\theta(\eta\kappa(\widehat\rho))$ is a nonlinear function of the estimated
correlation, so symmetric noise in $\widehat\rho$ shifts its mean as well as
spreading it, and the direction of the shift depends on the cluster.
Damping by $\eta$ shrinks the spread, which always helps, and moves the
mean, which helps or hurts according to that direction. The sign of
\eqref{eq:crit} is the outcome of the competition.

\subsection{Both dials}

Take two clusters of two assets. Each has a unit-variance knot $p$ and a
member $m$ of volatility $s$ correlated $r$ with its knot, the knots have
correlation $c$, and the members' cross-cluster covariance passes through
the knots, so that $\mathrm{Cov}(m_1,p_2)=r_1s_1c$ and
$\mathrm{Cov}(m_1,m_2)=r_1s_1r_2s_2c$. Let all three correlations be
estimated as $r_1\pm\tau$, $r_2\pm\tau$, $c\pm\tau$ independently. The
verification script computes the derivatives in \cref{prop:corner} as
truncated power series with rational coefficients, so \eqref{eq:shift} is
exact.

With $(r_1,r_2,c,s_1,s_2)=(\tfrac14,\tfrac12,\tfrac12,2,3)$ both shifts
are negative, about $-29.9\,\tau^2$ in $\gamma$ and $-11.8\,\tau^2$ in
$\eta$, so the optimum is strictly inside the square. With
$(\tfrac14,\tfrac12,\tfrac14,\tfrac12,3)$ the $\gamma$ shift is positive,
about $+15.1\,\tau^2$, and the $\eta$ shift negative, so the optimum sits
on the edge $\gamma=1$ at $\eta\approx1-16.6\,\tau^2$ by the edge formula.
At $\tau=\tfrac1{40}$ the predicted point beats the corner in both
examples in exact arithmetic. Noise in the cross-cluster correlation alone
still moves $\eta$, through the mixed term of $G$, so the two dials are not
separable in general.

\section{Cost}

At $\gamma=0$ the square costs one inverse of each cluster block, from
which \eqref{eq:closed} gives the whole inner dial and the fitness. The
outer dial costs the conditioning. With the knot device of
\cref{sec:square} nothing larger than a cluster or the number of clusters
is inverted at any point of the square, the budget of the bridge in
\cite{cotton2026nco}. HERC itself inverts nothing, and the edge $\gamma=0$
adds one inverse per cluster to it.

\section{Closing}

Until now HERC has been placed with the heuristic methods, as a variant of
hierarchical risk parity that budgets by risk contribution. Its budget rule
is the global minimum-variance rule. What separates it from the optimum is
conditioning, omitted twice, and each omission is a Schur complement that
can be restored by degrees. The dendrogram it builds is used, in the
precision reading of its own split formula, only to choose the clusters.

Two things are new in the inner dial. Stevens' identity becomes a
one-parameter path with a closed form, a long-only frontier, and, on an
equicorrelated cluster, a straight segment whose noise problem is one
scalar. And that scalar decides, in closed form, whether estimation error
should stop the path short of the cluster's minimum-variance portfolio.

Open directions are the analogues of those in \cite{cotton2026nco}: lost
precision and effective damping for the outer dial under a gateway model,
global statements for the joint problem, and the non-quadratic risk
measures HERC also admits, for which no fitness identity of the form
\eqref{eq:fitness} exists and the square has no exact far corner.

\begin{thebibliography}{99}

\bibitem{raffinot2018}
T. Raffinot.
The hierarchical equal risk contribution portfolio.
SSRN 3237540, 2018.

\bibitem{lopezdeprado2016}
M. L\'opez de Prado.
Building diversified portfolios that outperform out of sample.
\emph{Journal of Portfolio Management}, 42(4):59--69, 2016.

\bibitem{cotton2024schur}
P. Cotton.
Schur complementary allocation: a unification of hierarchical risk parity
and minimum variance portfolios.
arXiv:2411.05807, 2024.

\bibitem{cotton2026bridge}
P. Cotton.
When the out-of-sample-optimal Schur portfolio lies between HRP and
minimum variance.
Working paper, 2026.
\url{https://schur.microprediction.org/neither-end-of-the-bridge.pdf}

\bibitem{cotton2026nco}
P. Cotton.
Nested clustered optimization is one end of a Schur bridge, and the
interior is sometimes provably better.
SSRN 7480738, 2026.

\bibitem{stevens1998}
G. V. G. Stevens.
On the inverse of the covariance matrix in portfolio analysis.
\emph{Journal of Finance}, 53(5):1821--1827, 1998.

\bibitem{noguer2026}
M. Noguer i Alonso.
The mathematics of heuristic portfolio optimization.
arXiv:2606.12612, 2026.

\bibitem{palomar2025}
D. P. Palomar.
\emph{Portfolio Optimization: Theory and Application}.
Cambridge University Press, 2025. Section 12.3.

\bibitem{choueifaty2008}
Y. Choueifaty and Y. Coignard.
Toward maximum diversification.
\emph{Journal of Portfolio Management}, 35(1):40--51, 2008.

\bibitem{arevalo2019}
R. Ar\'evalo, J. Le\'on, and G. Hern\'andez.
Portfolio selection based on hierarchical clustering and inverse-variance
weighting.
In \emph{Computational Science, ICCS 2019}, LNCS 11538, 2019.

\bibitem{trucios2026}
C. Truc\'ios.
Hierarchical risk clustering versus traditional risk-based portfolios: an
empirical out-of-sample comparison.
\emph{Empirical Economics}, 2026.

\bibitem{huang2022}
Y. Huang and X. Gao.
Evaluating hierarchical equal risk contribution portfolios in the Chinese
stock market.
\emph{Journal of Mathematical Finance}, 12:179--195, 2022.

\bibitem{sjostrand2020}
D. Sj\"ostrand and N. Behnejad.
Exploration of hierarchical clustering in long-only risk-based portfolio
optimization.
Master's thesis, Copenhagen Business School, 2020.

\bibitem{riskfolio}
D. Cajas.
Riskfolio-Lib. \texttt{github.com/dcajasn/Riskfolio-Lib}, 2020--2026.

\bibitem{skfolio}
H. Delatte and C. Nicolini.
skfolio. \texttt{github.com/skfolio/skfolio}, 2023--2026.

\bibitem{mlfinlab}
Hudson and Thames.
PortfolioLab, formerly mlfinlab, 2019--2021.

\end{thebibliography}

\end{document}
